\section{Packing over a handle}\label{sec:hdl}

The handle construction turns the comparison inequality into ball embeddings by three exact
symplectic paths. It is an explicit version of the surface packing lemma of \cite[\S3.2]{OpenAI2026}, in which the
Moser argument \cite{Moser1965} is carried out with explicit vector fields. The first path puts the usable base mass
on an annulus, the second inserts the comparator isotopy across it, and the third separates the caps into disjoint
layers. The higher-dimensional models use the finite coordinate construction of \Ssec{sec:native-toric}; the
four-dimensional periodic model uses the equivariant version of Remark~\ref{rem:hdl-four}. The suspension and layer
identities below have the same form in both.

In the higher-dimensional construction below, $n\ge3$ and $m=n-1$. Write $\Omega=\omega_V+d\Gamma_{\rm phys}$ on the explicit affine product, with compact model region
$V_\Delta(E)=\Sigma_\ell\times V_\Delta$, physical base curvature $\kappa(p)$ and profile
$\Hb(p)=\int_{\Sigma_\ell}\kappa(p)$. The comparison datum has gap $\gamma$ and reserve
$0<\gamma_*\le\gamma$. A horizontal one-form annihilates fibre tangent vectors. A family of base
one-forms depending smoothly on $p$ gives a toric horizontal form; its differential is
$d_\Sigma\alpha(p)+\sum_jdp_j\wedge\partial_j\alpha(p)$. This expression is smooth across the
coordinate axes because the moments $p_j=\pi|w_j|^2$ are smooth there.

\subsection{Handle data}\label{sec:hdl-data}

\begin{definition}[Handle data]\label{def:hdl-handle-data}
\lean{SymplecticPacking.HighDim.RawAnnularHandle,
SymplecticPacking.HighDim.ToricPackingStages.quadricRawHandle,
SymplecticPacking.HighDim.ToricPackingStages.quadricModelRawPresentation}
Each higher-dimensional model supplies the following data.
\begin{enumerate}[label=(\alph*)]
\item An explicit annular coordinate map on $|s_{\rm a}|<w$, with working band
$|s_{\rm a}|\le b<w$. We use $s=(s_{\rm a}+b)/(2b)$, so the working band is $[0,1]$.
For the diagonal-quadric presentation the model uses $w=1/8$ and $b=1/16$ inside a double cylinder of half-width $1/4$ on which the annular coordinate is defined.
\item A fixed compact support set $\mathcal B\subset\operatorname{int}\Sigma_\ell$.
\item A smooth closed one-form $\beta$, supported in $\mathcal B$, whose annular expression is
$b(v)\,dv$, with $b\ge0$, $b>0$ on $(v_-,v_+)$ and $\int_0^1b(v)\,dv=1$.
\item An explicit reference primitive $\Gamma_{\rm ref}$ with positive curvature
$\Lambda=d_\Sigma\Gamma_{\rm ref}$, a reference-area upper bound $A_{\rm ref}>0$, and a positive
physical density floor $\kappa_-$ relative to $\Lambda$ on the correction region:
$\kappa(p)\ge\kappa_-\Lambda$ for $p\in\Delta$.
\item A finite coordinate-box presentation of that region: explicit chart maps and inverses,
partition weights, unit-mass coordinate bumps, and a rooted tree on the boxes in which every box other than the root has an overlapping parent of smaller index.
Rectangular integration along this tree defines the operator $\mathcal P$ of \Ssec{sec:native-toric}.
\item The product fibre coordinates of the affine model. In the four-dimensional instance these
are the explicit equivariant coordinates on the cover.
\item A collar cutoff $\chi_\partial$, equal to one near the boundary and zero on the inner region
containing the annulus, together with the model's proved tail bound. A second, larger finite atlas
covers the whole closed model region for the Moser flows.
\end{enumerate}
\end{definition}

Every item is supplied by the model formulas. The fixed annulus and dual form are constructed in
\Ssec{sec:vtr}; the finite atlas and tree are in \Ssec{sec:native-toric}.
The annular gauge is identically one on $|s_{\rm a}|<b_1=(b+w)/2$.
We write $\mathcal A^+$ for this open band and $\mathcal A$ for the smaller open band
$|s_{\rm a}|<(b+b_1)/2$. Both contain the working band $0\le s\le1$ in their interior.
Thus the clock $f(s)=S(s)$ is constant on collars of the boundary of $\mathcal A$.
Put $\eta=df\wedge\beta$ and write
$\Lambda=a_\Lambda(s,v)\,ds\wedge dv$, $a_\Lambda>0$.
For the periodic four-dimensional instance the background $e_\Lambda\Lambda$ is $e_h\nu$ (Remark~\ref{rem:hdl-four}).

\begin{lemma}[Top power]\label{lem:hdl-top-power}
\lean{SymplecticPacking.Forms.pf_splitForm,
SymplecticPacking.Handle.Cubic.OmegaPreG_split,
OAI.PackingSufficiencySupport.Hamiltonian.toric_horizontalCoupling_isInvertible,
OAI.PackingSufficiencySupport.Hamiltonian.oneCovectorForm_isInvertible}
\uses{def:mech-model, def:hdl-handle-data}
\begin{enumerate}[label=(\alph*)]
\item In local product coordinates let $\kappa'$ be a base two-form and $a'_1,\dots,a'_m$ base one-forms, with
coefficients depending on the base point and on $p$. Then
$(\omega_V+\kappa'+\sum_jdp_j\wedge a'_j)^n=n\,\omega_V^m\wedge\kappa'$.
\item On $E|_{\mathcal A^+}$, in the explicit product coordinates, let $\sigma_\Sigma=a_\Sigma(s,v)\,ds\wedge dv$ be positive, let $G(s,w)$ be smooth and
independent of $v$, and $e>0$. Then
$\big(\omega_V+e\sigma_\Sigma+d(G\beta)\big)^n=n\,(e a_\Sigma(s,v)+b(v)\,\partial_sG)\,\omega_V^m\wedge ds\wedge dv$.
\end{enumerate}
In particular the forms in (a) are nondegenerate where $\kappa'>0$; when closed, they are
symplectic there. The forms in (b) are symplectic where $\partial_sG\ge0$.
\end{lemma}

\begin{proof}
Write $x_j,y_j$ for the real fibre coordinate covectors and $s,t$ for the base
covectors at a point. Every $dp_j$ has the form $u_jx_j+v_jy_j$ and every $a'_j$
has the form $A_js+B_jt$. For $B=A_js+B_jt$, the identity
\[
 x_j\wedge y_j+(u_jx_j+v_jy_j)\wedge B
 =(x_j-v_jB)\wedge(y_j+u_jB)
\]
factors each block as a decomposable two-form. These blocks commute and square
to zero. Their sum therefore has zero $(m+1)$-st power, and its $m$-th power is
$m!$ times their product. Against the base area, every base correction in that
product vanishes. Adding $\kappa'$ now gives (a) by the binomial formula. This is a
Cartesian identity, including points at which a moment differential vanishes.

For (b), write $d(G\beta)=\partial_sG\,ds\wedge\beta+d_wG\wedge\beta$.
Its mixed part has rank one: it is a fibre covector wedged with the single base
covector $\beta$. Apply the same block factorization with a common base covector
to obtain $n\omega_V^m\wedge(e\sigma_\Sigma+\partial_sG\,ds\wedge\beta)$.
The displayed coefficient follows from $\beta=b(v)\,dv$. Positive coefficients
make the top power nonzero. Thus a closed form of either displayed type is symplectic.
\end{proof}

\subsection{The first path}\label{sec:hdl-first}

Set
\begin{equation}\label{eq:hdl-eh}
 e_\Lambda=\min\left(\frac{\gamma_*}{32A_{\rm ref}},\frac{\kappa_-}{2}\right)>0,
 \qquad \epsilon_0=\frac{\gamma_*}{16}.
\end{equation}
The model tail estimate and the reference-area bound give
\begin{equation}\label{eq:hdl-collar}
 \int_{\Sigma_\ell}\chi_\partial\kappa(p)\le\frac{\gamma_*}{32},\qquad
 e_\Lambda\int_{\Sigma_\ell}\Lambda\le\frac{\gamma_*}{32}\qquad(p\in\Delta).
\end{equation}
Apply the finite primitive operator $\mathcal P$ of \Ssec{sec:native-toric}, a finite sum of rectangular integrals over the boxes corrected along the tree so that $d\mathcal P\omega=\omega$ for every two-form $\omega$ of total mass zero, to
$(1-\chi_\partial)(\kappa(p)-e_\Lambda\Lambda)$, then subtract its annular period and an exact annular gauge.
The resulting primitive $\Gamma_n$ is the formula in \Sref{Construction}{cons:native-normalization}.
It has
\begin{equation}\label{eq:hdl-ksm}
\begin{gathered}
 \kappa_{\rm sm}(p)=\chi_\partial\kappa(p)+(1-\chi_\partial)e_\Lambda\Lambda,\qquad
 \epsilon_{\rm sm}(p)=\int_{\Sigma_\ell}\kappa_{\rm sm}(p),\\
 H^*(p)=\Hb(p)-\epsilon_{\rm sm}(p),\qquad
 d_\Sigma\Gamma_n(p)=\kappa_{\rm sm}(p)+H^*(p)\eta .
\end{gathered}
\end{equation}
The mass equality is proved by calibration with the root bump (\Sref{Lemma}{lem:native-calibration}); the
construction itself uses the ratio of two finite coordinate integrals. Since
$H^*(p)=\int_{\Sigma_\ell}(1-\chi_\partial)(\kappa(p)-e_\Lambda\Lambda)$, the support of $1-\chi_\partial$ lies in
the correction region, and $\kappa(p)\ge\kappa_-\Lambda\ge2e_\Lambda\Lambda$ there, we have $H^*\ge0$ on $\Delta$;
and by \eqref{eq:hdl-collar}
\begin{equation}\label{eq:hdl-budget}
 0\le\epsilon_{\rm sm}\le\epsilon_0\le\gamma/16,\qquad
 H^*\ge\Hb-\gamma_*/16.
\end{equation}
The background is positive, agrees with the physical curvature near the outer boundary, and equals
$e_\Lambda\Lambda$ on the annulus. There the period and gauge corrections give
$\Gamma_n=e_\Lambda\Gamma_{\rm ref}+f(s)H^*(p)\beta$ on an open neighbourhood of the working band.
Define
\begin{equation}\label{eq:hdl-alpha1}
 \alpha_1=\Gamma_n-\Gamma_{\rm phys}.
\end{equation}

\begin{proposition}[First path]\label{prop:hdl-first-path}
\lean{SymplecticPacking.HighDim.RawModelAgreement.normalization_target,
SymplecticPacking.HighDim.RawModelAgreement.relativePath_semantic,
SymplecticPacking.Bundle.alpha1_curvature, SymplecticPacking.Bundle.Cubic.alpha1_annulus}
The form $\Omega_\lambda=\Omega+\lambda d\alpha_1$, $0\le\lambda\le1$, is symplectic on a
neighbourhood of the model region. Its base curvature is the positive convex combination of
$\kappa$ and $d_\Sigma\Gamma_n$, and on the annulus its endpoint is
\[
 \Omega_1=\omega_V+e_\Lambda\Lambda+d\big(f(s)H^*(p)\beta\big).
\]
The variation primitive is horizontal and toric, is smooth in the moment variables, and vanishes
near the outer base boundary.
\end{proposition}
\begin{proof}
The normalisation identities give the base curvature and the annular normal form. A toric mixed term
does not affect the top power, so the positive convex combination is nondegenerate. The difference
$\Gamma_n-\Gamma_{\rm phys}$ has compact support in the interior of the base, which gives the boundary
statement. All moment derivatives are derivatives of the explicit primitive formulas.
\end{proof}

\subsection{The twist}\label{sec:hdl-twist}

The generators $Q_t$ satisfy $\partial_tq_t=X_{Q_t}\circ q_t$. On a collar of each outer face the $Q_t$
Poisson-commute with $\mu_\nu$, so $\mu_\nu$ is constant along the trajectories of the isotopy there. The inverse
isotopy $q_t^{-1}$ is generated by $-Q_t\circ q_t$, which Poisson-commutes with $\mu_\nu$ on a slightly smaller collar
because $q_t$ preserves $\mu_\nu$ there. Hence neither $q_t$ nor $q_t^{-1}$ moves a point across an outer face:
$q_t(V_\Delta)=V_\Delta$, $q_t(\operatorname{int}V_\Delta)=\operatorname{int}V_\Delta$, and $q_t$ maps each outer face into
itself. We write $q=q_1$.

Over the annulus, in the explicit product coordinates, the coefficient of $\beta$ in $\Omega_1$
rises from $0$ to $H^*(p)$. The twist changes the fibre identification across the annulus from the identity to $q$.
For $\lambda\in[0,1]$ define on $E|_{\mathcal A}$
\begin{equation}\label{eq:hdl-twist}
 \mathcal D_\lambda(s,v,w)=\big(s,v,q_{\lambda f(s)}(w)\big),\qquad
 Q^\lambda_s=\lambda f'(s)\,Q_{\lambda f(s)},\qquad
 G_\lambda(s,w)=f(s)\,H^*\big(p(q_\lambda^{-1}(w))\big),
\end{equation}
and $\mathcal D=\mathcal D_1$. Since $\partial_s\big(q_{\lambda f(s)}(w)\big)=\lambda f'(s)\,(\partial_tq_t)_{t=\lambda f(s)}(w)$, the
function $Q^\lambda_s$ generates the $s$-dependence: $\partial_sq_{\lambda f(s)}=X_{Q^\lambda_s}\circ q_{\lambda f(s)}$.

\begin{lemma}[Suspension]\label{lem:hdl-suspension}
\lean{SymplecticPacking.HighDim.RawModelAgreement.firstPath_semantic, SymplecticPacking.Handle.Cubic.suspension, SymplecticPacking.Handle.Cubic.eta_support}
\uses{def:hdl-handle-data, def:mech-datum, prop:hdl-first-path}
On $E|_{\mathcal A}$, for $\lambda\in[0,1]$:
\begin{enumerate}[label=(\roman*)]
\item $\mathcal D_\lambda^*\omega_V=\omega_V+d\big[(Q^\lambda_s\circ q_{\lambda f(s)})\,ds\big]$;
\item with the horizontal one-form
\[
 \eta_\lambda=(Q^\lambda_s\circ q_{\lambda f(s)})\,ds
 +f(s)\Big[H^*\big(p(q_\lambda^{-1}q_{\lambda f(s)}(w))\big)-H^*\big(p(w)\big)\Big]\beta ,
\]
one has
$\mathcal D_\lambda^*\big(\omega_V+e_\Lambda\Lambda+d(G_\lambda\beta)\big)=\omega_V+e_\Lambda\Lambda+d\big(f(s)H^*(p)\beta\big)+d\eta_\lambda$;
\item $\eta_0=0$, $\eta_\lambda$ is smooth in $(\lambda,s,v,w)$, and $\eta_\lambda=0$ for $s\notin(0,1)$.
Hence $\eta_\lambda$ extends by zero to a horizontal one-form on a neighbourhood of $V_\Delta(E)$, supported over
$\mathcal A$, and $\Omega_1+d\eta_\lambda=\mathcal D_\lambda^*\big(\omega_V+e_\Lambda\Lambda+d(G_\lambda\beta)\big)$ on $E|_{\mathcal A}$.
\end{enumerate}
\end{lemma}

\begin{proof}
(i) Put $\Theta(s,w)=Q^\lambda_s(q_{\lambda f(s)}(w))$, independent of $v$; the right side is $\omega_V+d_w\Theta\wedge ds$. The
differential of $\mathcal D_\lambda$ fixes $\partial_v$, maps a vertical vector $u$ to $Dq\,u$
($q=q_{\lambda f(s)}$) and $\partial_s$ to $\partial_s+X_{Q^\lambda_s}(q(w))$; the form $\omega_V$ sees only vertical
components. On two vertical vectors both sides give $\omega_V(u_1,u_2)$, because $q$ is symplectic. On
$(\partial_s,u)$ the left side is $\omega_V(X_{Q^\lambda_s},Dq\,u)=-dQ^\lambda_s(Dq\,u)=-d_w\Theta(u)$, which is also
$(d_w\Theta\wedge ds)(\partial_s,u)$. Pairs containing $\partial_v$, and $(\partial_s,\partial_v)$, give zero on both sides.
(ii) $\mathcal D_\lambda$ fixes base points, so it preserves $\Lambda$ and $\beta$ and
$\mathcal D_\lambda^*(G_\lambda\beta)=(G_\lambda\circ\mathcal D_\lambda)\beta$, with
$G_\lambda\circ\mathcal D_\lambda=f(s)H^*(p(q_\lambda^{-1}q_{\lambda f(s)}w))$; add (i).
(iii) For $s\le0$, $f=f'=0$ and both terms vanish. For $s\ge1$, $f=1$, $f'=0$ and
$q_\lambda^{-1}q_{\lambda f(s)}=q_\lambda^{-1}q_\lambda=\mathrm{id}$, so the bracket vanishes exactly. The last identity is (ii)
together with Proposition~\ref{prop:hdl-first-path}.
\end{proof}

Using the fixed map $q_\lambda^{-1}$ in $G_\lambda$, rather than $q_{\lambda f(s)}^{-1}$, makes the bracket vanish identically for
$s\ge1$ and gives $\partial_sG_\lambda=f'(s)H^*(p(q_\lambda^{-1}w))$, with no derivative of the isotopy.

\emph{The second path.} Put $\Omega^{(1)}_\lambda=\Omega_1+d\eta_\lambda$. Off $E|_{\mathcal A}$ it equals $\Omega_1$. On
$E|_{\mathcal A}$ it is $\mathcal D_\lambda^*$ of a form whose top power is, by Lemma~\ref{lem:hdl-top-power}(b),
$n(e_\Lambda a_\Lambda(s,v)+b\,f'H^*(p(q_\lambda^{-1}w)))\,\omega_V^m\wedge ds\wedge dv$. This is positive on a neighbourhood of $V_\Delta(E)$:
the background $e_\Lambda\Lambda$ is positive and $H^*\ge0$ at every fibre moment in $\Delta$. Smoothness and compactness then give an open nondegenerate neighbourhood (\Sref{Lemma}{lem:shd-nbhd}); $\mathcal D_\lambda$ preserves orientation. So $\Omega^{(1)}_\lambda$, $\lambda\in[0,1]$, is an
exact symplectic path (\Sref{Proposition}{prop:shd-second-path}). The map $\mathcal D_\lambda$ itself is defined only over
$\mathcal A$; only the forms are glued to $E$.

\subsection{Layers and the Moser map}\label{sec:hdl-layers}

Fix source capacities $\rho_i^\sharp<a_i$ and put $\sigma_i=(a_i-\rho_i^\sharp)/2>0$.
Define the smooth caps and the collar coefficient by
\begin{equation}\label{eq:hdl-caps}
 h_i(p)=(a_i-|p|)S\big((a_i-|p|)/\sigma_i\big)\quad(1\le i\le k),\qquad
 J=K-H^*=K-\Hb+\epsilon_{\rm sm} .
\end{equation}
Put $N=k+1$ and $D_s=1/(3N+1)$. For $j=1,\dots,N$ set
$l_j=(3j-2)D_s$, $r_j=(3j-1)D_s$, $I_j=[l_j,r_j]$ and
$\varsigma_j(s)=S\big(2(s-l_j-D_s/4)/D_s\big)$; the last index $j=N$ is denoted by $*$.
The intervals are ordered and disjoint in $(0,1)$, and each clock is constant on collars
of its interval endpoints, with $\varsigma_j'\ge0$.

The constants enter below only through \eqref{eq:hdl-budget} and the residual bound
$\gamma_{\rm res}=\gamma-\sup_\Delta\epsilon_{\rm sm}\ge15\gamma/16$.

\begin{proposition}[Layers]\label{prop:hdl-layers}
\lean{SymplecticPacking.HighDim.ToricPackingStages.smoothCap_smooth,
SymplecticPacking.HighDim.ToricPackingStages.smoothCap_nonneg,
SymplecticPacking.HighDim.ToricPackingStages.smoothCap_le,
SymplecticPacking.HighDim.ToricPackingStages.smoothCap_ge}
\uses{def:mech-datum, lem:hdl-suspension, S-lem:shd-nbhd}
Under the hypotheses of Theorem~\ref{thm:mech-packing}, with $\epsilon_{\rm sm}$ from \eqref{eq:hdl-ksm} and $h_i,J$ from \eqref{eq:hdl-caps}:
\begin{enumerate}[label=(\roman*)]
\item each $h_i$ is smooth on $\R^m$ ($2n\ge6$; for $n=2$ see Remark~\ref{rem:hdl-four}), and
$0\le h_i(p)\le(a_i-|p|)_+$ and $(a_i+\rho_i^\sharp)/2-|p|\le h_i(p)$;
it vanishes for $|p|\ge a_i$ and equals $a_i-|p|$ for $|p|\le(a_i+\rho_i^\sharp)/2$;
\item the residual $R^*(w)=K\big(p(q^{-1}w)\big)-J\big(p(w)\big)-\sum_{i=1}^kh_i\big(p(w)\big)$ satisfies, for $w\in V_\Delta$,
\[
 R^*(w)\ \ge\ \gamma-\sup_\Delta\epsilon_{\rm sm}\ \ge\ \frac{15\gamma}{16}>0 ;
\]
\item the final coefficient
$G_{\rm f}(s,w)=J(p(w))+\sum_{i=1}^k\varsigma_i(s)h_i(p(w))+\varsigma_*(s)R^*(w)$ satisfies  $\partial_sG_{\rm f}\ge0$ on
$\mathcal A\times V_\Delta$; $G_{\rm f}=J\circ p$ for $s\le l_1$; $G_{\rm f}=K\circ p\circ q^{-1}$ for
$s\ge r_*$; and on $I_i\times V_\Delta$, $1\le i\le k$,
\[
 G_{\rm f}(s,w)=A_i(p(w))+\varsigma_i(s)h_i(p(w)),\qquad A_i=J+\sum_{j<i}h_j .
\]
\end{enumerate}
\end{proposition}

\begin{proof}
(i) Write $x=a_i-|p|$. For $x\le0$ the cap is zero. For $0<x<\sigma_i$, the bounds
$0\le S(x/\sigma_i)\le1$ give $0\le xS(x/\sigma_i)\le x$ and $x-\sigma_i\le0\le h_i$.
For $x\ge \sigma_i$ the cap equals $x$. Thus $x-\sigma_i\le h_i\le x_+$ in all cases.
(ii) Let $w\in V_\Delta$ and $z=q^{-1}(w)\in V_\Delta$, so that $w=q(z)$. Using $J=K-\Hb+\epsilon_{\rm sm}$, the hypothesis
$\Hb-\sum_i(a_i-|p|)_+\ge P$ and (i),
\begin{align*}
R^*(w)&=K(p(z))-K(p(qz))+\Hb(p(qz))-\epsilon_{\rm sm}(p(qz))-\sum_ih_i(p(qz))\\
&\ge K(p(z))-(K-P)(p(qz))-\sup_\Delta\epsilon_{\rm sm}
\ \ge\ \gamma-\sup_\Delta\epsilon_{\rm sm},
\end{align*}
by Definition~\ref{def:mech-datum} at the point $z$. Finally $\epsilon_{\rm sm}\le\gamma/16$ by \eqref{eq:hdl-budget}.
(iii) $\partial_sG_{\rm f}=\sum_i\varsigma_i'h_i+\varsigma_*'R^*\ge0$ by (i), and because $R^*>0$ on $V_\Delta$ by (ii). For $s\le l_1$ all $\varsigma_j$ vanish; for $s\ge r_*$
all equal one and $J+\sum_ih_i+R^*=K\circ p\circ q^{-1}$. On $I_i$, $\varsigma_j=1$ for $j<i$ and $\varsigma_j=0$ for $j>i$,
including $j=*$.
\end{proof}

The comparison is used at the source point $z=q^{-1}(w)$: the twist transports the budget $K$ of level $p(z)$ to the
fibre point $w$ where it is spent. The residual layer comes last, so all non-toric dependence on $w$ is confined to
$I_*$ and the ball layers see only functions of $p$.

\emph{The third path.} Keep $\mathcal D$ fixed and interpolate the coefficient from $G_1=f(s)H^*\circ p\circ q^{-1}$ to
$G_{\rm f}$. On $E|_{\mathcal A}$ put $\xi_{\mathcal A}=((G_{\rm f}-G_1)\circ\mathcal D)\,\beta$. On the lower collar
$s\le0$, $\mathcal D=\mathrm{id}$, $G_1=0$ and $G_{\rm f}=J\circ p$; on the upper collar $s\ge1$,
$G_1\circ\mathcal D=H^*\circ p$ and $G_{\rm f}\circ\mathcal D=K\circ p\circ q^{-1}\circ q=K\circ p$. So
$\xi_{\mathcal A}=J(p)\beta$ on both collars, and
\begin{equation}\label{eq:hdl-xi}
 \xi=J(p)\,\mathrm{pr}^*\beta+\operatorname{ext}_0\big(\xi_{\mathcal A}-J(p)\beta\big),
\end{equation}
where $\operatorname{ext}_0$ denotes extension by zero (the form $\xi_{\mathcal A}-J(p)\beta$ vanishes over both collars of
$\mathcal A$), is a horizontal one-form on a neighbourhood of $V_\Delta(E)$; the first term is a global toric form,
since $p$ and $\beta$ are. Put $\Omega^{(2)}_\tau=\Omega^{(1)}_1+\tau\,d\xi$ and $\Omega_{\rm final}=\Omega^{(2)}_1$. On
$E|_{\mathcal A}$,
\begin{equation}\label{eq:hdl-final}
\begin{aligned}
 \Omega^{(2)}_\tau&=\mathcal D^*\big(\omega_V+e_\Lambda\Lambda+d(\hat G_\tau\beta)\big),&\hat G_\tau&=(1-\tau)G_1+\tau G_{\rm f},\\
 \Omega_{\rm final}&=\mathcal D^*\Omega_{\rm pre},&\Omega_{\rm pre}&=\omega_V+e_\Lambda\Lambda+d(G_{\rm f}\beta);
\end{aligned}
\end{equation}
off $E|_{\mathcal A}$, $\Omega^{(2)}_\tau=\Omega_1+\tau\,dJ\wedge\beta$ differs from $\Omega_1$ by a toric mixed term. By
Lemma~\ref{lem:hdl-top-power}, Proposition~\ref{prop:hdl-layers}(iii) and \Sref{Lemma}{lem:shd-nbhd}, every $\Omega^{(2)}_\tau$ is
symplectic on a neighbourhood of $V_\Delta(E)$ (\Sref{Proposition}{prop:shd-layer-path}). Two different toric coefficients at the two
ends of $\mathcal A$ would have to be interpolated across the rest of $\Sigma$, changing its area; the single form
$J(p)\mathrm{pr}^*\beta$ adds only the mixed term $dJ\wedge\beta$, which does not enter the top power.

For each of the three paths write its primitive as $\lambda_{j,t}=\lambda_V+Q_{j,t}$ and its variation
as $\nu_{j,t}=\partial_tQ_{j,t}$; these are $\alpha_1$, $\partial_t\eta_t$ and $\xi$, respectively.
In the larger finite atlas use the whole-path Gram cutoff $\chi_j$ of
\Sref{Construction}{cons:native-packing}. It is one near $[0,1]\times V_\Delta(E)$ and has compact spatial
support in the nondegenerate locus. Define the cut field by
$\iota_{Y_{j,t}}d\lambda_{j,t}=-\chi_j\nu_{j,t}$, computed in the charts of the finite atlas, and let $\Phi_j$ be
its time-one map. Put $\Phi=\Phi_2\circ\Phi_1\circ\Phi_0$ and use the abbreviation
$\Phi^{-1}=\Phi_0^{-1}\circ\Phi_1^{-1}\circ\Phi_2^{-1}$.

\begin{proposition}[Moser maps]\label{prop:hdl-faces}
\lean{SymplecticPacking.HighDim.RawFlowAtlasSpec.pathFlow_eq_horizontalMoserMap,
SymplecticPacking.HighDim.ToricPackingStages.RawSurfacePresentation.fullMoserFlow_diffeomorph,
SymplecticPacking.HighDim.RawModelAgreement.native_packing_spec}
The three cut fields are smooth with common compact spatial support. Their time-one maps are
diffeomorphisms, preserve the model region and its interior in both directions, and satisfy
$\Phi_j^*\Omega_{j,1}=\Omega_{j,0}$ there. Consequently
\[
 \Phi^*\Omega_{\rm final}=\Omega.
\]
\end{proposition}
\begin{proof}
Near each outer face the variation primitives annihilate the Hamiltonian field of the face moment. Every path form is
invariant under the flow of that moment there: the toric terms are, and the comparator isotopy commutes with it
because $Q_t$ Poisson-commutes with $\mu_\nu$ on the collar (Definition~\ref{def:mech-datum}\ref{it:mech-datum-collar}),
so $\iota_{X_{\mu_\nu}}\Omega_{j,t}=-d\mu_\nu$ on the collar.
The comparator's face-collar condition therefore gives $d\mu_\nu(Y_{j,t})=0$ whenever that moment is
near its boundary value. A trajectory cannot cross the corresponding level in either direction;
this argument applies to all faces, including their intersections. The variation primitives vanish
near the outer base boundary. The fields have compact support, so their flows are complete and their
inverse maps are the backward flows. On the model region $\chi_j=1$, and Cartan's identity gives
\[
 \frac d{dt}\Phi_{j,t}^*\Omega_{j,t}
 =\Phi_{j,t}^*d\bigl(\iota_{Y_{j,t}}\Omega_{j,t}+\nu_{j,t}\bigr)=0.
\]
The finite coordinate ODEs and their inverse endpoints are exactly the maps used in the construction.
\end{proof}

\emph{Layer coordinates.} On the layer $\Lambda_i=\{l_i<s<r_i,\ v_-<v<v_+\}\times\operatorname{int}V_\Delta$ (product
coordinates) let $T(v)=\int_{v_-}^vb$, a diffeomorphism $(v_-,v_+)\to(0,1)$ with $dT=\beta$, and
\begin{equation}\label{eq:hdl-layer-x}
 \ell(s,T)=\frac{e_\Lambda a_\Lambda(s,v(T))}{b(v(T))},\qquad
 x_i(s,T,p)=\int_{l_i}^s\ell(u,T)\,du+\varsigma_i(s)h_i(p).
\end{equation}
Then $\partial_sx_i=\ell+\varsigma_i'h_i>0$, $x_i=0$ at $s=l_i$ and $x_i>h_i(p)$ at $s=r_i$, so for $0<x<h_i(p)$ and
$0<T<1$ there is a unique root $s=s_i(x,T,p)\in(l_i,r_i)$, smooth by the implicit function theorem. Since
$e_\Lambda\Lambda=\ell(s,T)\,ds\wedge dT$ and, by Proposition~\ref{prop:hdl-layers}(iii),
$d(G_{\rm f}\beta)=(\varsigma_i'h_i\,ds+dA_i+\varsigma_i\,dh_i)\wedge dT$ on $\Lambda_i$,
\begin{equation}\label{eq:hdl-layer-form}
 \Omega_{\rm pre}=\omega_V+dx_i\wedge dT+dA_i(p)\wedge dT\qquad\text{on }\Lambda_i .
\end{equation}

Two explicit maps finish the layer. The \emph{multi-angle shear} $\mathrm{Sh}_A$, which rotates each fibre coordinate
$w_j$ by the angle $-2\pi T\partial_jA(p(w))$, removes the term $dA_i(p)\wedge dT$ from~\eqref{eq:hdl-layer-form}
(\Sref{Lemma}{lem:hdl-shear}). The \emph{disc-to-strip map} $\Psi_\delta\colon\C\to(0,\infty)\times(0,1)$ is an explicit
symplectomorphism onto its image with $0<x<\pi|z|^2+\delta/2$ (\Sref{Proposition}{prop:hdl-strip}); the excess $\delta/2$
is unavoidable, since a closed disc of area $a$ does not embed area-preservingly in $\{0<x<a\}\times(0,1)$, and it is paid
once for all concentric discs.

\subsection{Proof of Theorem~\ref{thm:mech-packing}}\label{sec:hdl-proof}

For a higher-dimensional input, let the model, handle data, capacities $a_i$, function $P$ and comparison datum $(K,q)$ with gap $\gamma$ be as in
Theorem~\ref{thm:mech-packing}, with the bounds \eqref{eq:hdl-collar}, and let
$\rho_i<a_i$. Write $\C^n=\C\times\C^m\ni(z_{\rm b},z_{\rm f})$, $\varpi=\pi|z_{\rm b}|^2$, $p=p(z_{\rm f})$, so that
$\omega_0=\omega_\C(z_{\rm b})+\omega_V(z_{\rm f})$ and $B(c)=\{\varpi+|p|\le c\}$.

\emph{Parameters.} Put
\[
 \hat\rho_i=\frac{3a_i+\rho_i}4,\qquad \varepsilon_i=\frac{a_i-\hat\rho_i}4,
\]
so that $\rho_i<\hat\rho_i<a_i$ and $B(\rho_i)\subset\Bo(\hat\rho_i)$; the open ball $\Bo(\hat\rho_i)$ is the domain of
the embedding. In \eqref{eq:hdl-caps} take $\rho_i^\sharp=\hat\rho_i$; these are the constants of
\Sref{equation}{eq:shd-constants}.

\emph{Forms and Moser map.} Construct $c$, $\kappa_{\rm sm}$, $\alpha_1$, the paths $\Omega_\lambda$, $\Omega^{(1)}_\lambda$,
$\Omega^{(2)}_\tau$, the twist $\mathcal D$, the coefficient $G_{\rm f}$ and the Moser map $\Phi$ as above. By
Proposition~\ref{prop:hdl-faces}, $\Phi^*\Omega_{\rm final}=\Omega$, and by \eqref{eq:hdl-final}
$\Omega_{\rm final}=\mathcal D^*\Omega_{\rm pre}$ on $E|_{\mathcal A}$, where $\mathcal D$ is a diffeomorphism of
$\mathcal A\times V_\Delta$ fixing base points, with inverse $(s,v,w)\mapsto(s,v,q_{f(s)}^{-1}(w))$ preserving
$\operatorname{int}V_\Delta$.

\emph{Layer maps.} For $(z_{\rm b},z_{\rm f})\in B^\circ(\hat\rho_i)$ let $(x,T)=\Psi_{\varepsilon_i}(z_{\rm b})$
(\Sref{Proposition}{prop:hdl-strip}). Then $|p|<\hat\rho_i<a_i$, so $z_{\rm f}\in\operatorname{int}V_\Delta$ because
$\{|p|\le a_i\}\subset\tau_\Delta\Delta$, and
\[
 0<x<\varpi+\frac{\varepsilon_i}2
 <\hat\rho_i-|p|+\frac{a_i-\hat\rho_i}{8}
 <\frac{a_i+\hat\rho_i}{2}-|p|\le h_i(p)
\]
by Proposition~\ref{prop:hdl-layers}(i). Define
\[
 \mathcal E_i(z_{\rm b},z_{\rm f})=\Big(s_i\big(x,T,p\big),\ v(T),\ \mathrm{Sh}_{A_i,T}(z_{\rm f})\Big)\in\Lambda_i ,
\]
with the root $s_i$ of \eqref{eq:hdl-layer-x} and the shear of \Sref{Lemma}{lem:hdl-shear} ($\mathrm{Sh}_{A_i,T}$ preserves $p$).
$\mathcal E_i$ is the composite of $\Psi_{\varepsilon_i}\times\mathrm{id}$, the shear $\mathrm{Sh}_{A_i}$ and the inverse of the layer
coordinates $(s,v,w)\mapsto(x_i,T,w)$, each a diffeomorphism onto an open set; by \eqref{eq:hdl-layer-form},
\Sref{Lemma}{lem:hdl-shear} and \Sref{Proposition}{prop:hdl-strip},
\[
 \mathcal E_i^*\Omega_{\rm pre}=(\Psi_{\varepsilon_i}\times\mathrm{id})^*\big(\omega_V+dx\wedge dT\big)=\omega_\C+\omega_V=\omega_0 .
\]
Put
\[
 \psi_i=\Phi_0^{-1}\circ\Phi_1^{-1}\circ\Phi_2^{-1}
             \circ\mathcal D^{-1}\circ\mathcal E_i
       =\Phi^{-1}\circ\mathcal D^{-1}\circ\mathcal E_i
       \colon B^\circ(\hat\rho_i)\to E .
\]

In \Sref{Construction}{cons:native-packing} the map $\mathcal D^{-1}\circ\mathcal E_i$ is written as a composite whose factors are the planar map, the unique inverse layer chart, the inverse angular clock,
the annular covering map and the inverse comparator shear. The endpoint factors are, in order,
the inverse flows of $Q_0=(1-t)\Gamma_{\rm phys}+t\Gamma_n$,
$Q_1=\Gamma_n+\eta_t$ and $Q_2=\Gamma_n+\eta_1+t\xi$.
Thus $\psi_i$ is the map of \Sref{Construction}{cons:native-packing} with $\rho_i^\sharp=\hat\rho_i=(3a_i+\rho_i)/4$.

\emph{Conclusions.} (a) $\psi_i^*\Omega=\mathcal E_i^*(\mathcal D^{-1})^*(\Phi^{-1})^*\Omega=\mathcal E_i^*(\mathcal D^{-1})^*\Omega_{\rm final}
=\mathcal E_i^*\Omega_{\rm pre}=\omega_0$, and $\psi_i$ is an injective immersion of an open set, hence a smooth symplectic embedding of
the neighbourhood $\Bo(\hat\rho_i)$ of $B(\rho_i)$. (b) The image of $\mathcal E_i$ lies in
$\{l_i<s<r_i\}\times\operatorname{int}V_\Delta$, and $\mathcal D^{-1}$, $\Phi^{-1}$ map $\operatorname{int}V_\Delta(E)$ into itself
(Proposition~\ref{prop:hdl-faces}), so $\psi_i$ takes values in $\operatorname{int}V_\Delta(E)$. (c) The intervals
$(l_i,r_i)$ are pairwise disjoint, so the images of the $\mathcal E_i$ are pairwise disjoint.
The maps $\mathcal D^{-1}$ and $\Phi^{-1}$ are common injective maps for all $i$, and therefore preserve this disjointness.


Every ingredient is given by a formula, a fixed linear operator, a definite integral, the unique solution of an explicit ODE
($\Phi$) or a unique root ($s_i$), so $\psi_i$ is explicit in the sense of Definition~\ref{def:intro-levels}.
The verification of these maps is \Sref{Theorem}{thm:shd-packing}, with the strip and
layer estimates in \Sref{Proposition}{prop:str-layer} and \Sref{Proposition}{prop:str-embeddings}. \qed

The source margin $a_i-\rho_i$ pays for the strip excess and for an open neighbourhood of each given
closed ball. For $2n\ge6$ the caps lie below the original tents and their deficit is retained in the final residual
layer, while the background and collar consume at most $\gamma_*/16$; in dimension four the caps lie above the tents
by a total of at most $g/32$, which the residual budget absorbs (\Sref{Lemma}{lem:shd4-budget}).

\begin{remark}[Dimension four]\label{rem:hdl-four}
For $n=2$ the base is the closed torus of \Sref{Section}{sec:cub}, with the handle data of
\Sref{Proposition}{prop:cub-handle}; there is no collar and no atlas, and the construction is carried out on the cover
$\R^2\times\C$, equivariantly under the deck maps (\Ssec{sec:shd-four}). The background is $e_h\nu$ with $e_h=g/16$,
$\nu=du\wedge dv$, and the first path is $\Omega_0+t\,d\alpha_1$ with the closed-form primitive $\alpha_1$ of
\Sref{Construction}{cons:fer-alpha1}. The twist uses the clock $f(u)=S((u-\frac7{48})/\frac1{24})$, which is
$S(2s-\frac12)$ in the handle coordinate $s=12u-\frac32$, and the $k+1$ layer intervals lie inside its ramp. The caps are $h_i(p)=\frac12\bigl(a_i-p+\sqrt{(a_i-p)^2+\mu_h^2}\bigr)$ with
$\mu_h=g/(16k)$; they lie above the tents and exceed them in total by at most $g/32$, so a comparison gap
$\frac{31}{32}g$ leaves a residual of at least $\frac78g$. The three stages are concatenated into one flow, cut off
only in the moment variable, on $\{p<h+2\delta_m\}$, $\delta_m=\min((d_n-h)/4,g/100)$. All three variation
primitives vanish as covectors on the marked fibres, so the flow fixes them pointwise and the images avoid
$\{Z_0=0\}$ (\Sref{Lemma}{lem:shd4-fibres}, \Sref{Propositions}{prop:shd4-moser} and~\Snum{prop:shd4-maps}).
\end{remark}
